% Options for packages loaded elsewhere
\PassOptionsToPackage{unicode}{hyperref}
\PassOptionsToPackage{hyphens}{url}
\documentclass[
]{article}
\usepackage{xcolor}
\usepackage{amsmath,amssymb}
\setcounter{secnumdepth}{-\maxdimen} % remove section numbering
\usepackage{iftex}
\ifPDFTeX
  \usepackage[T1]{fontenc}
  \usepackage[utf8]{inputenc}
  \usepackage{textcomp} % provide euro and other symbols
\else % if luatex or xetex
  \usepackage{unicode-math} % this also loads fontspec
  \defaultfontfeatures{Scale=MatchLowercase}
  \defaultfontfeatures[\rmfamily]{Ligatures=TeX,Scale=1}
\fi
\usepackage{lmodern}
\ifPDFTeX\else
  % xetex/luatex font selection
\fi
% Use upquote if available, for straight quotes in verbatim environments
\IfFileExists{upquote.sty}{\usepackage{upquote}}{}
\IfFileExists{microtype.sty}{% use microtype if available
  \usepackage[]{microtype}
  \UseMicrotypeSet[protrusion]{basicmath} % disable protrusion for tt fonts
}{}
\makeatletter
\@ifundefined{KOMAClassName}{% if non-KOMA class
  \IfFileExists{parskip.sty}{%
    \usepackage{parskip}
  }{% else
    \setlength{\parindent}{0pt}
    \setlength{\parskip}{6pt plus 2pt minus 1pt}}
}{% if KOMA class
  \KOMAoptions{parskip=half}}
\makeatother
\usepackage{color}
\usepackage{fancyvrb}
\newcommand{\VerbBar}{|}
\newcommand{\VERB}{\Verb[commandchars=\\\{\}]}
\DefineVerbatimEnvironment{Highlighting}{Verbatim}{commandchars=\\\{\}}
% Add ',fontsize=\small' for more characters per line
\newenvironment{Shaded}{}{}
\newcommand{\AlertTok}[1]{\textcolor[rgb]{1.00,0.00,0.00}{\textbf{#1}}}
\newcommand{\AnnotationTok}[1]{\textcolor[rgb]{0.38,0.63,0.69}{\textbf{\textit{#1}}}}
\newcommand{\AttributeTok}[1]{\textcolor[rgb]{0.49,0.56,0.16}{#1}}
\newcommand{\BaseNTok}[1]{\textcolor[rgb]{0.25,0.63,0.44}{#1}}
\newcommand{\BuiltInTok}[1]{\textcolor[rgb]{0.00,0.50,0.00}{#1}}
\newcommand{\CharTok}[1]{\textcolor[rgb]{0.25,0.44,0.63}{#1}}
\newcommand{\CommentTok}[1]{\textcolor[rgb]{0.38,0.63,0.69}{\textit{#1}}}
\newcommand{\CommentVarTok}[1]{\textcolor[rgb]{0.38,0.63,0.69}{\textbf{\textit{#1}}}}
\newcommand{\ConstantTok}[1]{\textcolor[rgb]{0.53,0.00,0.00}{#1}}
\newcommand{\ControlFlowTok}[1]{\textcolor[rgb]{0.00,0.44,0.13}{\textbf{#1}}}
\newcommand{\DataTypeTok}[1]{\textcolor[rgb]{0.56,0.13,0.00}{#1}}
\newcommand{\DecValTok}[1]{\textcolor[rgb]{0.25,0.63,0.44}{#1}}
\newcommand{\DocumentationTok}[1]{\textcolor[rgb]{0.73,0.13,0.13}{\textit{#1}}}
\newcommand{\ErrorTok}[1]{\textcolor[rgb]{1.00,0.00,0.00}{\textbf{#1}}}
\newcommand{\ExtensionTok}[1]{#1}
\newcommand{\FloatTok}[1]{\textcolor[rgb]{0.25,0.63,0.44}{#1}}
\newcommand{\FunctionTok}[1]{\textcolor[rgb]{0.02,0.16,0.49}{#1}}
\newcommand{\ImportTok}[1]{\textcolor[rgb]{0.00,0.50,0.00}{\textbf{#1}}}
\newcommand{\InformationTok}[1]{\textcolor[rgb]{0.38,0.63,0.69}{\textbf{\textit{#1}}}}
\newcommand{\KeywordTok}[1]{\textcolor[rgb]{0.00,0.44,0.13}{\textbf{#1}}}
\newcommand{\NormalTok}[1]{#1}
\newcommand{\OperatorTok}[1]{\textcolor[rgb]{0.40,0.40,0.40}{#1}}
\newcommand{\OtherTok}[1]{\textcolor[rgb]{0.00,0.44,0.13}{#1}}
\newcommand{\PreprocessorTok}[1]{\textcolor[rgb]{0.74,0.48,0.00}{#1}}
\newcommand{\RegionMarkerTok}[1]{#1}
\newcommand{\SpecialCharTok}[1]{\textcolor[rgb]{0.25,0.44,0.63}{#1}}
\newcommand{\SpecialStringTok}[1]{\textcolor[rgb]{0.73,0.40,0.53}{#1}}
\newcommand{\StringTok}[1]{\textcolor[rgb]{0.25,0.44,0.63}{#1}}
\newcommand{\VariableTok}[1]{\textcolor[rgb]{0.10,0.09,0.49}{#1}}
\newcommand{\VerbatimStringTok}[1]{\textcolor[rgb]{0.25,0.44,0.63}{#1}}
\newcommand{\WarningTok}[1]{\textcolor[rgb]{0.38,0.63,0.69}{\textbf{\textit{#1}}}}
\usepackage{longtable,booktabs,array}
\newcounter{none} % for unnumbered tables
\usepackage{calc} % for calculating minipage widths
% Correct order of tables after \paragraph or \subparagraph
\usepackage{etoolbox}
\makeatletter
\patchcmd\longtable{\par}{\if@noskipsec\mbox{}\fi\par}{}{}
\makeatother
% Allow footnotes in longtable head/foot
\IfFileExists{footnotehyper.sty}{\usepackage{footnotehyper}}{\usepackage{footnote}}
\makesavenoteenv{longtable}
\setlength{\emergencystretch}{3em} % prevent overfull lines
\providecommand{\tightlist}{%
  \setlength{\itemsep}{0pt}\setlength{\parskip}{0pt}}
\usepackage{fancyhdr}
\usepackage{eso-pic}
\usepackage{tikz}
\usepackage{geometry}
\usepackage{xcolor}
\geometry{margin=1in}
\AddToShipoutPictureBG{%
  \begin{tikzpicture}[remember picture, overlay]
    \node[rotate=45, scale=3.5, text opacity=0.10, color=gray]
      at (current page.center)
      {\textsf{Gem v0.1.0 \textperiodcentered\ 2026-04-18}};
  \end{tikzpicture}%
}
\pagestyle{fancy}
\fancyhf{}
\fancyhead[L]{\small\textbf{Gem Language} v0.1.0}
\fancyhead[R]{\small 2026-04-18}
\fancyfoot[C]{\small\thepage}
\renewcommand{\headrulewidth}{0.4pt}
\renewcommand{\footrulewidth}{0.4pt}
\usepackage{bookmark}
\IfFileExists{xurl.sty}{\usepackage{xurl}}{} % add URL line breaks if available
\urlstyle{same}
\hypersetup{
  hidelinks,
  pdfcreator={LaTeX via pandoc}}

\title{Gem Language Tutorial Series}
\author{David B Hon}
\date{2026-04-18}

\begin{document}
\maketitle

\section{David B Hon Hard Science Gem Language Copywrite 2025,
2026}\label{david-b-hon-hard-science-gem-language-copywrite-2025-2026}

\textbf{Version:} 0.1.0 \textbar{} \textbf{Date:} 2026-04-18 \textbar{}
\textbf{Built with:} Gemini CLI \& Kiro CLI

Welcome to the \textbf{Gem Language} --- a modern, expressive STEM
language built by Gemini-CLI and Kiro-CLI using C++26 for AI, scientific
computing, quantitative finance, and cross-platform application
development.

\subsection{Language Mindmap}\label{language-mindmap}

🔍 Interactive mindmap --- drag to pan, scroll to zoom

\subsection{Key Features}\label{key-features}

\begin{itemize}
\tightlist
\item
  \textbf{Explicit Typing with Initial Values}: All variable
  declarations require an initial value (\texttt{int\ x\ =\ 0}).
\item
  \textbf{Global vs.~Local Scope}: Variables starting with \texttt{\_}
  (e.g., \texttt{\_config}) are global; all others are local.
\item
  \textbf{Universal Inheritance}: Every object inherits from
  \texttt{sys}, giving immediate access to \texttt{sys.print}, etc.
\item
  \textbf{First-Class AI Support}: Built-in \texttt{ai} object supports
  Gemini, Mistral, and Ollama.
\item
  \textbf{Polyglot Interop}: \texttt{use} keyword bridges Python, Julia,
  R, Fortran, C++, Go, Ruby, Rust, Node.js on-the-fly. Native execution
  via \texttt{python.run()}, \texttt{python.compile()},
  \texttt{cython.build()}, \texttt{go.run()}, etc.
\item
  \textbf{Mobile \& Cross-Platform}: \texttt{mobl} object + browser PWA
  (Web Speech + Geolocation) works on Android, iPhone, macOS, Linux, and
  Windows 11.
\item
  \textbf{Domain-Specific Built-ins}: \texttt{fin}, \texttt{bsm},
  \texttt{bev}, \texttt{geo}, \texttt{data}, \texttt{chart},
  \texttt{mobl}, \texttt{trek}, \texttt{astro}, \texttt{nlp},
  \texttt{www}, \texttt{cdn}, \texttt{thread}, \texttt{rex},
  \texttt{seo}, \texttt{drvr}, \texttt{art}, \texttt{git} built into the
  runtime (37 builtins total: includes \texttt{python} and
  \texttt{cython}).
\item
  \textbf{Regular Expressions}: \texttt{rex} builtin provides full
  ECMAScript regex --- match, find, findall, groups, sub, gsub, split,
  count.
\item
  \textbf{Symbolic Math}: \texttt{math} builtin supports symbolic
  differentiation, integration, simplification, and LaTeX output via
  SymPy/Sage.
\item
  \textbf{Astrophysics}: \texttt{astro} builtin covers stellar physics,
  orbital mechanics, cosmology, solar physics, and exoplanets.
\end{itemize}

\begin{center}\rule{0.5\linewidth}{0.5pt}\end{center}

\subsection{Tutorial Roadmap}\label{tutorial-roadmap}

\subsubsection{1. The Foundations}\label{the-foundations}

\begin{itemize}
\tightlist
\item
  \textbf{\url{01_basics.g}}: Variables, types, and basic I/O using the
  \texttt{sys} object.
\item
  \textbf{\url{02_math_latex.g}}: Arithmetic operations and LaTeX-style
  mathematical expressions.
\item
  \textbf{\url{03_text_pdf.g}}: Text manipulation and PDF document
  generation.
\item
  \textbf{\url{04_images.g}}: Image processing with ImageMagick.
\item
  \textbf{\url{05_geo_gis.g}}: GIS, geolocation, GeoJSON, and
  interactive maps.
\end{itemize}

\subsubsection{2. Application
Development}\label{application-development}

\begin{itemize}
\tightlist
\item
  \textbf{\url{06_web_apps.g}}: Building web servers with
  \texttt{sys.app()} and route handlers.
\item
  \textbf{\url{07_ai_nlp.g}}: NLP and AI prompting with Gemini, Mistral,
  and Ollama.
\item
  \textbf{\url{08_functions.gm}}: Defining functions, objects, and
  inheritance.
\item
  \textbf{\url{09_module_demo.g}}: Structuring code with modules.
\item
  \textbf{\url{11_modules_paths.g}}: Managing module resolution and
  import paths.
\end{itemize}

\subsubsection{3. Advanced Integration \&
AI}\label{advanced-integration-ai}

\begin{itemize}
\tightlist
\item
  \textbf{\url{10_interop_cpp.g}}: High-performance C++ interoperability
  via Cling JIT.
\item
  \textbf{\url{12_lib_management.g}}: Dependency management and library
  linking.
\item
  \textbf{\url{13_dependency_details.g}}: Deep dive into Gem's package
  ecosystem.
\item
  \textbf{\url{14_ai_providers.g}}: Using Gemini, Mistral, and Ollama
  side-by-side.
\item
  \textbf{\url{15_mistral_native.g}}: Native Mistral C++ bridge via
  \texttt{prompt\_native}.
\end{itemize}

\subsubsection{4. Specialized Tooling}\label{specialized-tooling}

\begin{itemize}
\tightlist
\item
  \textbf{\url{16_algo_text_extras.g}}: Advanced string algorithms and
  text analysis.
\item
  \textbf{\url{17_repl_features.g}}: Mastering the interactive Gem REPL.
\item
  \textbf{\url{18_tcp_sockets.g}}: Low-level networking fundamentals.
\item
  \textbf{\url{19_tcp_server.g}} \& \textbf{\url{20_tcp_client.g}}:
  Client-server architectures.
\item
  \textbf{\url{21_polyglot.g}}: Bridging Python, Julia, and Gem.
\end{itemize}

\subsubsection{5. Science, Finance \&
DevOps}\label{science-finance-devops}

\begin{itemize}
\tightlist
\item
  \textbf{\url{22_builtins_itr_tcp.g}}: Iterators and network protocols.
\item
  \textbf{\url{23_data_science.g}}: DataFrames, statistics, and
  visualization.
\item
  \textbf{\url{24_devops.g}}: Automation scripts and infrastructure
  management.
\item
  \textbf{\url{25_ollama_mistral.g}}: Running local LLMs via Ollama.
\item
  \textbf{\url{26_devops_containers_vms.g}}: Docker and Vagrant
  environments.
\end{itemize}

\subsubsection{6. High-Performance \& Symbolic
Math}\label{high-performance-symbolic-math}

\begin{itemize}
\tightlist
\item
  \textbf{\url{27_rust_node.g}}: Polyglot support for Rust and Node.js.
\item
  \textbf{\url{28_interactive_charts.g}}: Dynamic 2D/3D/multi-axes
  Plotly visualization.
\item
  \textbf{\url{29_quantlib_finance.g}}: Quantitative finance with
  QuantLib.
\item
  \textbf{\url{30_bsm_pde_inheritance.g}}: Black-Scholes-Merton PDE
  solvers.
\item
  \textbf{\url{31_history_and_compilation.g}}: Gem's compilation model
  and history.
\item
  \textbf{\url{32_symbolic_math.g}}: Symbolic differentiation and
  integration.
\item
  \textbf{\url{33_langport_porting.g}}: Porting legacy codebases to Gem.
\item
  \textbf{\url{34_geo_tectonics.g}}: Advanced GIS, tectonic modeling,
  and geo-visualization.
\item
  \textbf{\url{35_text_advanced.g}}: Large-scale text processing and
  indexing.
\end{itemize}

\subsubsection{7. Mobile \& Cross-Platform}\label{mobile-cross-platform}

\begin{itemize}
\tightlist
\item
  \textbf{\url{36_mobl_travel_log.g}}: \texttt{mobl} object --- NLP
  dictation + GPS + GeoJSON travel log PWA.
\end{itemize}

\subsubsection{8. System Programming \&
Hardware}\label{system-programming-hardware}

\begin{itemize}
\tightlist
\item
  \textbf{\url{37_device_drivers.g}}: \texttt{drvr} builtin --- Linux,
  Windows 11, macOS, and Android device driver development.
\end{itemize}

\subsubsection{9. Text \& Pattern Matching}\label{text-pattern-matching}

\begin{itemize}
\tightlist
\item
  \textbf{\url{38_rex.g}}: \texttt{rex} builtin --- full ECMAScript
  regular expressions: match, find, findall, groups, sub, gsub, split,
  count.
\end{itemize}

\subsubsection{10. Travel Logs \& OSM
Mapping}\label{travel-logs-osm-mapping}

\begin{itemize}
\tightlist
\item
  \textbf{\url{39_trek.g}}: \texttt{trek} builtin --- create, display,
  and edit travel logs on OpenStreetMap with GeoJSON/GPX support.
\end{itemize}

\subsubsection{11. Search Engine
Optimization}\label{search-engine-optimization}

\begin{itemize}
\tightlist
\item
  \textbf{\url{40_seo.g}}: \texttt{seo} builtin --- crawl URLs, extract
  SEO signals (title, meta, headings, links), write JSON index, and
  print analysis report.
\end{itemize}

\subsubsection{12. ASCII Art \& SVG}\label{ascii-art-svg}

\begin{itemize}
\tightlist
\item
  \textbf{\url{41_art.g}}: \texttt{art} builtin --- generate ASCII art
  from text, convert to/from SVG, render mindmap/README/tutorial as
  ASCII art.
\end{itemize}

\subsubsection{13. Python \& Cython
Polyglot}\label{python-cython-polyglot}

\begin{itemize}
\tightlist
\item
  \textbf{\url{42_python_cython.g}}: \texttt{python} and \texttt{cython}
  builtins --- run/compile Python scripts, byte-compile \texttt{.py} →
  \texttt{.pyc}, transpile \texttt{.pyx} → C → \texttt{.so} native
  extension, pip package management. \#\#\# 14. Heterogeneous Arrays \&
  Dictionaries
\item
  \textbf{\url{43_arrays_dicts.g}}: \texttt{algo.array} and
  \texttt{algo.dict} --- allocate and manipulate heterogeneous
  multidimensional arrays of arbitrary content, and dictionary objects
  with arbitrary key/value tuples. Covers 1-D arrays, N-D matrices, dict
  basics, dict-of-arrays, and array-of-dicts table patterns.
\end{itemize}

\subsubsection{15. Git Version Control}\label{git-version-control}

\begin{itemize}
\tightlist
\item
  \textbf{\url{44_git.g}}: \texttt{git} builtin --- clone repositories,
  pull, commit, push, and the \texttt{gitsync} one-shot workflow (pull +
  commit + push + show summary).
\end{itemize}

\begin{center}\rule{0.5\linewidth}{0.5pt}\end{center}

\subsection{Help \& Builtins}\label{help-builtins}

\begin{Shaded}
\begin{Highlighting}[]
\NormalTok{{-}{-}{-} Gem Language Help {-}{-}{-}}
\NormalTok{Gem is a high{-}performance STEM language implemented in C++26.}

\NormalTok{Usage Tips:}
\NormalTok{  {-} All variables must be initialized: \textquotesingle{}int x = 0\textquotesingle{}}
\NormalTok{  {-} Global variables start with underscore: \textquotesingle{}\_config = 1\textquotesingle{}}
\NormalTok{  {-} Use \textquotesingle{}use "file.py"\textquotesingle{} to translate and run foreign code.}
\NormalTok{  {-} Use \textquotesingle{}langport("*.py", "out.gm")\textquotesingle{} to port and save code.}
\NormalTok{  {-} Use \textquotesingle{}lib\textquotesingle{} to see all loaded modules.}
\NormalTok{  {-} Use \textquotesingle{}his\textquotesingle{} to view today\textquotesingle{}s session history.}

\NormalTok{CLI Options:}
\NormalTok{  gem \textless{}file.g\textgreater{}              {-} Run a script.}
\NormalTok{  gem {-}c \textless{}main.g\textgreater{} [mod...]  {-} Compile multiple files into a binary.}
\NormalTok{  gem {-}o \textless{}name\textgreater{}             {-} Specify output name for compiler.}
\NormalTok{  gem {-}h                    {-} Print history and exit.}
\NormalTok{  gem {-}t \textless{}file\textgreater{} [{-}o output] {-} AI{-}assisted translation to Gem.}

\NormalTok{Available Builtin Modules:}
\NormalTok{  sys, math, ai, text, rex, algo, bev, file, zip, nlp, img, www, cdn, geo,}
\NormalTok{  cpp, tcp, itr, thread, data, k3s, vm, go, ruby, node, rust,}
\NormalTok{  python, cython,}
\NormalTok{  fin, bsm, chart, astro, mobl, trek, seo, drvr, art,}
\NormalTok{  git}

\NormalTok{Keywords for Documentation:}
\NormalTok{  fun, obj, use, alias, his, lib, end, if, else, while,}
\NormalTok{  int, double, string, bool, langport,}
\NormalTok{  true, false, null, nil, nan}
\NormalTok{  (exit/quit are identifiers, not keywords; \textquotesingle{}return\textquotesingle{} is not a keyword — use \textquotesingle{}end \textless{}expr\textgreater{}\textquotesingle{})}

\NormalTok{Builtin Functions:}
\NormalTok{  isnil(x)  isnan(x)  tonum(x)  tostr(x)  len(x)  type(x)}

\NormalTok{Operators:}
\NormalTok{  Arithmetic:   +  {-}  *  /}
\NormalTok{  Compound:     +=  {-}=  *=  /=}
\NormalTok{  Comparison:   ==  !=  \textgreater{}  \textgreater{}=  \textless{}  \textless{}=}
\NormalTok{  Logical:      !  (prefix not)  \&\&  ||}
\NormalTok{  Separator:    ;  (multiple statements per line)}
\NormalTok{  String concat: +}

\NormalTok{Mobile \& Cross{-}Platform:}
\NormalTok{  mobl phone("name")  phone.dictate(text)  phone.make\_feature(lat,lon,text)}
\NormalTok{  ./gem travel\_log.g  →  http://localhost:8080  (Android/iPhone/macOS/Linux/Win11)}

\NormalTok{Detailed help: help "topic"  or  help(topic)}
\NormalTok{Full reference: helpfull  or  sys.helpfull()  (writes helpfull\_DATETIME.md, opens in viewer)}
\NormalTok{Paged reference: helpless  or  sys.helpless()  (writes helpfull\_DATETIME.md, pipes through less)}
\end{Highlighting}
\end{Shaded}

\begin{center}\rule{0.5\linewidth}{0.5pt}\end{center}

\subsection{Builtin Functions \&
Literals}\label{builtin-functions-literals}

{\def\LTcaptype{none} % do not increment counter
\begin{longtable}[]{@{}lll@{}}
\toprule\noalign{}
Item & Type & Description \\
\midrule\noalign{}
\endhead
\bottomrule\noalign{}
\endlastfoot
\texttt{true} & literal bool & Boolean true \\
\texttt{false} & literal bool & Boolean false \\
\texttt{null} / \texttt{nil} & literal void & Null/void value
(interchangeable) \\
\texttt{nan} & literal double & IEEE 754 Not-a-Number \\
\texttt{isnil(x)} & function & \texttt{→\ bool} --- true if x is
nil/null/void \\
\texttt{isnan(x)} & function & \texttt{→\ bool} --- true if x is NaN \\
\texttt{tonum(x)} & function & \texttt{→\ double} --- parse to number;
nan on failure \\
\texttt{tostr(x)} & function & \texttt{→\ string} --- convert to
string \\
\end{longtable}
}

\subsection{Builtin Module Reference}\label{builtin-module-reference}

\subsubsection{sys --- System Root Module}\label{sys-system-root-module}

Every object inherits from \texttt{sys}. Methods: -
\texttt{sys.print(...)} --- print to stdout - \texttt{sys.args()} ---
CLI args as vector -
\texttt{sys.async(func,\ {[}params{]},\ {[}timeout{]})} --- background
thread; returns thread handle - \texttt{sys.sethandler(signals,\ func)}
--- custom signal handler - \texttt{sys.host()} --- OS info string -
\texttt{sys.exec(cmd)} --- run shell command -
\texttt{sys.doc({[}path{]})} --- read file comments -
\texttt{sys.help({[}topic{]})} --- interactive help (with cross-language
comparisons) - \texttt{sys.exit()} --- exit interpreter -
\texttt{sys.langport(pattern,\ {[}output{]})} --- AI-port foreign code
to Gem - \texttt{sys.redirect(url,\ {[}port{]})} --- HTTP redirect
server - \texttt{sys.app(port,\ {[}routes{]})} --- background HTTP
server

Comparisons: C++=\texttt{std::cout/std::thread/signal()} ·
Python=\texttt{print/threading/signal} · Julia=\texttt{println/@async} ·
Go=\texttt{fmt.Println/goroutine} · Ruby=\texttt{puts/Thread.new} ·
Rust=\texttt{println!/std::thread::spawn}

\subsubsection{math --- Mathematics \& Symbolic
Math}\label{math-mathematics-symbolic-math}

Standard: \texttt{sin(x)}, \texttt{cos(x)}, \texttt{sqrt(x)},
\texttt{math.pi}

Symbolic (SymPy/Sage backend): - \texttt{math.useSymPy()},
\texttt{math.useSage()} --- switch backend -
\texttt{math.diff(expr,\ {[}var{]})} --- symbolic derivative -
\texttt{math.integrate(expr,\ {[}var{]})} --- symbolic integral -
\texttt{math.simplify(expr)}, \texttt{math.solve(expr,\ {[}var{]})} -
\texttt{math.sym\_latex(expr)}, \texttt{math.to\_latex(val)},
\texttt{math.write\_latex(path,\ content)} -
\texttt{math.read\_latex(path)}, \texttt{math.parse\_latex(text)},
\texttt{math.compile\_latex(path)}

Comparisons: C++=\texttt{\textless{}cmath\textgreater{}/GiNaC} ·
Python=\texttt{math/sympy} · Julia=\texttt{sin()/Symbolics.jl} ·
Go=\texttt{math.Sin()} · Ruby=\texttt{Math::sin()} ·
Rust=\texttt{f64::sin()}

\subsubsection{ai --- AI Integration}\label{ai-ai-integration}

Supports Gemini (default), Mistral, Ollama: - \texttt{ai.prompt(text)}
--- send prompt to current provider - \texttt{ai.prompt\_native(text)}
--- Mistral C++ bridge (requires \texttt{mistralai} \textgreater= 1.0,
Python 3.14) - \texttt{ai.code(spec)} --- polyglot code generation →
writes \texttt{YYYYMMDD\_solution.\{py,go,cpp,rb,rs,js,jl\}}, returns
list of filenames - \texttt{ai.useMistral(model)},
\texttt{ai.useOllama(model,\ {[}host{]})}, \texttt{ai.useGemini()} -
\texttt{ai.key(key)}, \texttt{ai.setHost(host)},
\texttt{ai.setPath(path)} - Properties: \texttt{provider},
\texttt{model}, \texttt{host}

Comparisons: C++=\texttt{libcurl+json} ·
Python=\texttt{google-generativeai/mistralai} · Julia=\texttt{HTTP.jl} ·
Go=\texttt{generative-ai-go} · Ruby=\texttt{ruby-openai} ·
Rust=\texttt{async-openai}

\subsubsection{text --- String \& Document
Processing}\label{text-string-document-processing}

\begin{itemize}
\tightlist
\item
  \texttt{text.read(path)}, \texttt{text.sub(s,\ start,\ {[}len{]})},
  \texttt{text.replace(s,\ old,\ new)}
\item
  \texttt{text.write\_pdf(path,\ text)},
  \texttt{text.write\_pdf\_a(path,\ text)},
  \texttt{text.read\_pdf(path)}
\item
  \texttt{text.read\_markdown/write\_markdown},
  \texttt{text.read\_yaml/write\_yaml}
\item
  \texttt{text.read\_html/write\_html},
  \texttt{text.read\_xml/write\_xml}
\item
  \texttt{text.read\_fits\_header(path)},
  \texttt{text.read\_hdf\_header(path)}
\end{itemize}

Comparisons: C++=\texttt{std::string/fstream/pugixml} ·
Python=\texttt{str/pypdf/pyyaml/lxml} · Julia=\texttt{String/YAML.jl} ·
Go=\texttt{strings/encoding/xml} · Ruby=\texttt{String/prawn/nokogiri} ·
Rust=\texttt{String/lopdf/quick-xml}

\subsubsection{geo --- GIS \& Geolocation}\label{geo-gis-geolocation}

\begin{itemize}
\tightlist
\item
  \texttt{geo.lookup()} --- IP geolocation (\texttt{.lat},
  \texttt{.lon}, \texttt{.city}, \texttt{.country})
\item
  \texttt{geo.distance(lat1,\ lon1,\ lat2,\ lon2)} --- Haversine
  distance in km
\item
  \texttt{geo.write\_geojson(path,\ features)},
  \texttt{geo.history(plate)} --- tectonic history (AI)
\item
  \texttt{geo.plot2d(data,\ {[}layout{]})} --- 2D Plotly/OSM map
\item
  \texttt{geo.plot3d(data,\ {[}layout{]})} --- 3D globe
\end{itemize}

Comparisons: C++=\texttt{GDAL/GEOS} ·
Python=\texttt{geopandas/folium/geopy} · Julia=\texttt{GeoStats.jl} ·
Go=\texttt{paulmach/orb} · Ruby=\texttt{rgeo} · Rust=\texttt{geo\ crate}

\subsubsection{astro --- Astrophysics \& Planetary
Science}\label{astro-astrophysics-planetary-science}

Constants: \texttt{G}, \texttt{c}, \texttt{AU}, \texttt{pc},
\texttt{ly}, \texttt{Msun}, \texttt{Rsun}, \texttt{Lsun}, \texttt{Tsun},
\texttt{Mearth}, \texttt{Rearth}, \texttt{H0}, \texttt{sigma\_sb}

Unit conversions: \texttt{to\_ly(pc)}, \texttt{to\_pc(ly)},
\texttt{to\_au(m)}, \texttt{deg\_to\_rad(d)}, \texttt{rad\_to\_deg(r)}

Stellar: \texttt{luminosity(R,T)}, \texttt{stefan\_boltzmann(T,R)},
\texttt{wien(T)}, \texttt{abs\_magnitude(m,d\_pc)},
\texttt{distance\_modulus(d\_pc)}, \texttt{spectral\_class(T)},
\texttt{schwarzschild\_radius(M)}, \texttt{escape\_velocity(M,R)}

Orbital: \texttt{orbital\_period(a\_au,M\_msun)},
\texttt{orbital\_velocity(a\_m,M\_kg)}, \texttt{hill\_sphere(a,mp,ms)},
\texttt{roche\_limit(R,rho\_p,rho\_s)}, \texttt{synodic\_period(p1,p2)},
\texttt{planet(name)} → Mercury\ldots Neptune

Solar: \texttt{solar\_flux(dist\_au)},
\texttt{solar\_wind\_pressure(n\_cm3,v\_km\_s)},
\texttt{sunspot\_cycle(year)},
\texttt{parker\_spiral\_angle(v\_sw,dist\_au)},
\texttt{solar\_activity()}

Cosmology: \texttt{hubble\_distance(z)}, \texttt{redshift\_velocity(z)},
\texttt{lookback\_time(z)}, \texttt{critical\_density()}

Coordinates: \texttt{equatorial\_to\_galactic(ra,dec)},
\texttt{angular\_separation(ra1,dec1,ra2,dec2)}

Exoplanet: \texttt{transit\_depth(rp\_rearth,rs\_rsun)},
\texttt{habitable\_zone(L\_lsun)}, \texttt{equilibrium\_temp(L,d,A)}

Pulsar: \texttt{pulsar\_spindown(P\_s,\ Pdot)} →
\texttt{\{age\_yr,\ Bfield\_G,\ edot\_W\}}

Comparisons: C++=\texttt{SOFA/ERFA} · Python=\texttt{astropy} ·
Julia=\texttt{AstroLib.jl} · Go=\texttt{soniakeys/meeus} ·
Ruby=\texttt{astronomy\ gem} · Rust=\texttt{astro\ crate}

\subsubsection{seo --- Search Engine
Optimization}\label{seo-search-engine-optimization}

\begin{itemize}
\tightlist
\item
  \texttt{seo.index(urls,\ output\_path)} --- crawl list of URLs,
  extract SEO signals, write JSON index
\item
  \texttt{seo.analyze(index\_path)} --- load index JSON and print SEO
  report to stdout
\end{itemize}

Signals extracted: \texttt{title}, \texttt{description},
\texttt{keywords}, \texttt{og:title}, \texttt{og:description},
\texttt{canonical}, \texttt{word\_count}, \texttt{img\_count},
\texttt{img\_with\_alt}, \texttt{internal\_links},
\texttt{external\_links}, \texttt{h1}/\texttt{h2}/\texttt{h3} headings

Comparisons: C++=\texttt{libcurl+custom\ parser} ·
Python=\texttt{scrapy/beautifulsoup4} · Julia=\texttt{HTTP.jl/Gumbo.jl}
· Go=\texttt{colly/goquery} · Ruby=\texttt{nokogiri/mechanize} ·
Rust=\texttt{scraper/reqwest}

\subsubsection{art --- ASCII Art \& SVG}\label{art-ascii-art-svg}

\begin{itemize}
\tightlist
\item
  \texttt{art.text\_to\_art(text,\ {[}font{]})} --- render text as ASCII
  art via figlet/toilet (fallback: box)
\item
  \texttt{art.art\_to\_file(text,\ path)} --- write ASCII art string to
  a text file
\item
  \texttt{art.art\_to\_svg(text,\ {[}path{]})} --- convert ASCII art to
  SVG; returns SVG string
\item
  \texttt{art.svg\_to\_art(svg\_path,\ {[}width{]})} --- convert SVG to
  ASCII art via rsvg-convert + jp2a
\item
  \texttt{art.mindmap({[}path{]})} --- render
  \texttt{docs/gem\_mindmap.md} as ASCII art (mermaid mindmap)
\item
  \texttt{art.readme({[}path{]})} --- render \texttt{README.md} headings
  as ASCII art summary
\item
  \texttt{art.tutorial({[}path{]})} --- render
  \texttt{tutorial/README.md} as ASCII art summary
\end{itemize}

Optional deps: \texttt{figlet}/\texttt{toilet} (text→art) ·
\texttt{mmdc} (mermaid CLI) · \texttt{rsvg-convert}+\texttt{jp2a}
(SVG→art)

Comparisons: C++=\texttt{libaa/figlet\ CLI} ·
Python=\texttt{pyfiglet/ascii\_magic} · Julia=\texttt{no\ built-in} ·
Go=\texttt{no\ built-in} · Ruby=\texttt{artii\ gem} ·
Rust=\texttt{figlet-rs\ crate}

\subsubsection{trek --- Travel Logs \& OSM
Mapping}\label{trek-travel-logs-osm-mapping}

\begin{itemize}
\tightlist
\item
  \texttt{trek.new(path)} --- create empty GeoJSON travel log
\item
  \texttt{trek.add(path,\ lat,\ lon,\ {[}title{]},\ {[}note{]})} ---
  append waypoint
\item
  \texttt{trek.edit(path,\ index,\ title,\ note)} --- update waypoint at
  index
\item
  \texttt{trek.remove(path,\ index)} --- remove waypoint at index
\item
  \texttt{trek.load(path)} --- return GeoJSON string
\item
  \texttt{trek.show(path,\ {[}port{]})} --- serve interactive
  OSM/Leaflet map with live edit UI (default port 8090)
\item
  \texttt{trek.export\_gpx(geojson\_path,\ gpx\_path)} --- export to GPX
\item
  \texttt{trek.stats(path)} --- returns
  \texttt{\{.waypoints,\ .distance\_km\}}
\end{itemize}

Comparisons: C++=\texttt{GDAL/libgpx} · Python=\texttt{gpxpy/folium} ·
Julia=\texttt{GPX.jl} · Go=\texttt{tkrajina/gpxgo} ·
Ruby=\texttt{gpx\ gem} · Rust=\texttt{gpx\ crate}

\subsubsection{mobl --- Mobile \& Cross-Platform
PWA}\label{mobl-mobile-cross-platform-pwa}

\begin{verbatim}
mobl phone("device_name")
phone.dictate(spoken_text)          # NLP parse → JSON {title, note, tags}
phone.make_feature(lat, lon, text)  # GPS + dictation → GeoJSON Feature
\end{verbatim}

Browser supplies GPS (Geolocation API) + speech (Web Speech API). Gem
server handles NLP, GeoJSON assembly, and file persistence.

Routes: \texttt{/} → PWA HTML, \texttt{/log} → POST
\texttt{\{lat,lon,text\}}, \texttt{/data} → GET GeoJSON
FeatureCollection

Comparisons: C++=\texttt{Qt} · Python=\texttt{Flask+SpeechRecognition} ·
Julia=\texttt{no\ built-in} · Go=\texttt{net/http+custom\ JS} ·
Ruby=\texttt{Rails+Web\ Speech} · Rust=\texttt{actix-web+custom\ JS}

\subsubsection{fin --- Financial
Engineering}\label{fin-financial-engineering}

\begin{itemize}
\tightlist
\item
  \texttt{fin.ticker(symbol)} --- real-time market data via yfinance →
  \texttt{\{.price,.volume,.open,.high,.low,.close,.pe\_ratio,.market\_cap,.dividend\_yield\}}
\item
  \texttt{fin.high\_yield\_bonds()}, \texttt{fin.high\_yield\_etfs()},
  \texttt{fin.high\_yield\_equities()} --- top 50 from TradingView
\item
  \texttt{fin.dashboard()} --- serve financial dashboard on port 8082
\item
  \texttt{fin.bs\_price(type,\ strike,\ spot,\ rate,\ vol,\ t)} ---
  European option pricing → double
\item
  \texttt{fin.greeks(type,\ strike,\ spot,\ rate,\ vol,\ t)} →
  \texttt{\{.npv,\ .delta,\ .gamma,\ .theta,\ .vega\}}
\item
  \texttt{fin.date(d,m,y)} →
  \texttt{\{.day,.month,.year,.serial,.str\}},
  \texttt{fin.calendar(name)}, \texttt{fin.is\_holiday(cal,d,m,y)→bool},
  \texttt{fin.add\_days(d,m,y,n)→\{.day,.month,.year,.str\}},
  \texttt{fin.diff\_days(d1,m1,y1,d2,m2,y2)→double}
\end{itemize}

Comparisons: C++=\texttt{QuantLib} ·
Python=\texttt{yfinance/quantlib-python} ·
Julia=\texttt{MarketData.jl/QuantLib.jl} · Go=\texttt{piquette/quantlib}
· Ruby=\texttt{quantlib-ruby} · Rust=\texttt{ta\ crate}

\subsubsection{bsm ---
Black-Scholes-Merton}\label{bsm-black-scholes-merton}

Inherits all \texttt{fin} methods, plus: -
\texttt{bsm.price\_american(symbol,\ type,\ duration)} --- PDE-enhanced
American option pricing
(\texttt{weekly}/\texttt{monthly}/\texttt{quarterly})

Comparisons: C++=\texttt{QuantLib\ FDBlackScholesVanillaEngine} ·
Python=\texttt{quantlib-python\ AmericanOption} ·
Julia=\texttt{QuantLib.jl\ AmericanExercise}

\subsubsection{chart --- Interactive Plotting
(Plotly.js)}\label{chart-interactive-plotting-plotly.js}

\begin{itemize}
\tightlist
\item
  \texttt{chart.plot(traces,\ {[}layout{]})} --- generate interactive
  HTML chart
\item
  \texttt{chart.show(path)} --- open chart in browser
\item
  \texttt{chart.server(path,\ {[}port{]})} --- serve chart via
  background HTTP server
\end{itemize}

Comparisons: C++=\texttt{matplotlib-cpp/ROOT} ·
Python=\texttt{plotly/matplotlib/seaborn} ·
Julia=\texttt{Plots.jl/Makie.jl} · Go=\texttt{gonum/plot} ·
Ruby=\texttt{gruff} · Rust=\texttt{plotters\ crate}

\subsubsection{data --- Data Science}\label{data-data-science}

\begin{itemize}
\tightlist
\item
  \texttt{data.read\_csv(path)} --- read CSV into vector of row objects
\item
  \texttt{data.mean(vector)} --- arithmetic mean
\item
  \texttt{data.std(vector)} --- sample standard deviation
\end{itemize}

Comparisons: C++=\texttt{Eigen/Armadillo} · Python=\texttt{pandas/numpy}
· Julia=\texttt{DataFrames.jl/CSV.jl} · Go=\texttt{gonum/stat} ·
Ruby=\texttt{daru\ gem} · Rust=\texttt{polars\ crate}

\subsubsection{rex --- Regular Expressions
(ECMAScript)}\label{rex-regular-expressions-ecmascript}

All functions accept optional flags string (\texttt{"i"} =
case-insensitive): - \texttt{rex.match(text,\ pattern,\ {[}flags{]})} →
bool - \texttt{rex.find(text,\ pattern,\ {[}flags{]})} → first match
string - \texttt{rex.findall(text,\ pattern,\ {[}flags{]})} → list of
all matches - \texttt{rex.groups(text,\ pattern,\ {[}flags{]})} →
capture groups from first match -
\texttt{rex.sub(text,\ pattern,\ repl,\ {[}flags{]})} → replace first
match - \texttt{rex.gsub(text,\ pattern,\ repl,\ {[}flags{]})} → replace
all matches - \texttt{rex.split(text,\ pattern,\ {[}flags{]})} → split
on pattern - \texttt{rex.count(text,\ pattern,\ {[}flags{]})} → count
non-overlapping matches

Comparisons: C++=\texttt{std::regex} · Python=\texttt{re} ·
Julia=\texttt{Regex/match()} · Go=\texttt{regexp} ·
Ruby=\texttt{=\textasciitilde{}/gsub} · Rust=\texttt{regex\ crate}

\subsubsection{algo --- Algorithms, Arrays \&
Dictionaries}\label{algo-algorithms-arrays-dictionaries}

\begin{itemize}
\tightlist
\item
  \texttt{algo.add(...)} --- sum all numeric arguments
\item
  \texttt{algo.quicksort(v)},
  \texttt{algo.sort(v,\ {[}start{]},\ {[}end{]})} --- sort numeric
  vectors
\item
  \texttt{algo.now()} --- current local time as string
\item
  \texttt{algo.date\_add(ts,\ days)}, \texttt{algo.date\_diff(t1,\ t2)}
  --- date arithmetic
\item
  \texttt{algo.array({[}d1,d2,...{]},\ {[}fill{]})} --- allocate
  heterogeneous N-D array; fill default nil
\item
  \texttt{algo.array\_get(arr,\ i0,\ i1,\ ...)} --- index into nested
  array
\item
  \texttt{algo.array\_set(arr,\ val,\ i0,\ i1,\ ...)} --- set element
  in-place
\item
  \texttt{algo.array\_push(arr,\ val)} --- append to 1-D array
\item
  \texttt{algo.array\_len(arr)} --- outermost dimension length
\item
  \texttt{algo.array\_shape(arr)} --- dimensions as numeric vector
\item
  \texttt{algo.dict({[}k,v,...{]})} --- create dictionary, optional
  initial key/value pairs
\item
  \texttt{algo.dict\_set(d,\ key,\ val)} --- set key (mutates in-place)
\item
  \texttt{algo.dict\_get(d,\ key)} --- get value by key (nil if missing)
\item
  \texttt{algo.dict\_del(d,\ key)} --- remove key
\item
  \texttt{algo.dict\_has(d,\ key)} --- existence check
\item
  \texttt{algo.dict\_keys(d)} --- all keys as string vector
\item
  \texttt{algo.dict\_vals(d)} --- all values as heterogeneous array
\item
  \texttt{algo.dict\_len(d)} --- entry count
\end{itemize}

Comparisons:
C++=\texttt{std::sort/std::chrono/std::vector\textless{}std::any\textgreater{}/std::map}
· Python=\texttt{sorted()/datetime/list/dict} ·
Julia=\texttt{sort!()/Dates/Any{[}{]}/Dict} ·
Go=\texttt{sort.Slice/time/{[}{]}interface\{\}/map} ·
Ruby=\texttt{Array\#sort!/Time/Array/Hash} ·
Rust=\texttt{slice.sort/SystemTime/Vec\textless{}Box\textless{}dyn\ Any\textgreater{}\textgreater{}/HashMap}

\subsubsection{tcp --- TCP/IP Networking}\label{tcp-tcpip-networking}

\begin{itemize}
\tightlist
\item
  \texttt{tcp.listen(port)→int} --- bind server socket
\item
  \texttt{tcp.accept(fd)→\{.fd,.addr\}} --- block for client connection
\item
  \texttt{tcp.connect(host,\ port)→int} --- connect to server
\item
  \texttt{tcp.send(fd,\ msg)→int} --- send string, returns bytes sent
\item
  \texttt{tcp.recv(fd,\ {[}size{]})→string} --- receive data (default
  size 4096)
\item
  \texttt{tcp.close(fd)→bool} --- close socket
\item
  \texttt{tcp.nic()→{[}\{.name,.ip,.status\}{]}} --- network interface
  list
\item
  \texttt{tcp.routes()→{[}\{.dest,.gw,.iface\}{]}} --- routing table
\end{itemize}

Comparisons: C++=\texttt{POSIX\ socket/bind/listen} ·
Python=\texttt{socket.socket} · Julia=\texttt{Sockets.listen} ·
Go=\texttt{net.Listen} · Ruby=\texttt{TCPServer} ·
Rust=\texttt{TcpListener::bind}

\subsubsection{thread --- Background
Threads}\label{thread-background-threads}

Returned by \texttt{sys.async()}. Methods on the thread handle: -
\texttt{thread.wait()} --- block until complete, returns result (or
timeout error string) - \texttt{thread.is\_finished()} --- bool,
non-blocking status check

Comparisons: C++=\texttt{std::future/std::async} ·
Python=\texttt{concurrent.futures.Future} ·
Julia=\texttt{@async/fetch()} · Go=\texttt{goroutine+channel} ·
Ruby=\texttt{Thread\#join/alive?} · Rust=\texttt{JoinHandle\#join}

\subsubsection{www --- Web Framework}\label{www-web-framework}

\begin{itemize}
\tightlist
\item
  \texttt{www.wget(url,\ file)} --- download file via curl
\item
  \texttt{www.app(port,\ {[}routes{]})} --- background HTTP server;
  routes map paths to HTML strings
\item
  \texttt{www.redirect(target,\ {[}port{]})} --- start redirect server
  or return redirect header string
\item
  \texttt{www.map2d(geojson,\ output)} --- render 2D map via Mapnik
  (requires HAS\_MAPNIK)
\end{itemize}

Comparisons: C++=\texttt{cpp-httplib/Crow} ·
Python=\texttt{Flask/FastAPI} · Julia=\texttt{HTTP.jl/Genie.jl} ·
Go=\texttt{net/http/Gin} · Ruby=\texttt{Sinatra/Rails} ·
Rust=\texttt{actix-web/axum}

\subsubsection{cdn --- Caching Proxy
Server}\label{cdn-caching-proxy-server}

Inherits all \texttt{www} methods, plus: -
\texttt{cdn.start(port,\ routes\_map)} --- start caching reverse-proxy;
routes map path prefixes to origin URLs - \texttt{cdn.stats()} ---
returns object with \texttt{.hits}, \texttt{.misses}, \texttt{.bytes},
\texttt{.cached\_items} - \texttt{cdn.purge({[}path{]})} --- evict one
path or \texttt{"*"} to clear entire cache -
\texttt{cdn.config(type,\ content)} --- generate nginx/apache config -
\texttt{cdn.dashboard()} --- serve CDN stats dashboard on port 8083

Comparisons: C++=\texttt{nginx/Varnish} · Python=\texttt{no\ built-in} ·
Julia=\texttt{no\ built-in} · Go=\texttt{groupcache} ·
Ruby=\texttt{rack-cache} · Rust=\texttt{no\ built-in}

\subsubsection{k3s --- Docker \&
Kubernetes}\label{k3s-docker-kubernetes}

\begin{itemize}
\tightlist
\item
  \texttt{k3s.docker\_run(image,\ cmd)}, \texttt{k3s.docker\_ps()},
  \texttt{k3s.docker\_build(path,\ tag)}, \texttt{k3s.docker\_stop(id)}
\item
  \texttt{k3s.k3s\_apply(yaml)}, \texttt{k3s.k3s\_get(resource)},
  \texttt{k3s.k3s\_pods()}, \texttt{k3s.k3s\_nodes()},
  \texttt{k3s.k3s\_logs(pod)}
\end{itemize}

Comparisons: C++=\texttt{Docker\ SDK/HTTP\ API} ·
Python=\texttt{docker-py/kubernetes} · Julia=\texttt{no\ official\ SDK}
· Go=\texttt{docker/client/k8s.io/client-go} ·
Ruby=\texttt{docker-api/kubeclient} · Rust=\texttt{bollard/kube-rs}

\subsubsection{vm --- Vagrant VMs}\label{vm-vagrant-vms}

\begin{itemize}
\tightlist
\item
  \texttt{vm.init(box)}, \texttt{vm.up()}, \texttt{vm.ssh(cmd)},
  \texttt{vm.status()}, \texttt{vm.halt()}, \texttt{vm.destroy()}
\end{itemize}

Comparisons: C++=\texttt{libvirt\ C\ API} ·
Python=\texttt{python-vagrant} · Julia=\texttt{no\ official\ bindings} ·
Go=\texttt{os/exec} · Ruby=\texttt{Vagrant\ is\ Ruby} ·
Rust=\texttt{std::process::Command}

\subsubsection{drvr --- Device Driver
Development}\label{drvr-device-driver-development}

\begin{itemize}
\tightlist
\item
  \texttt{drvr.linux(name)}, \texttt{drvr.win11(name)},
  \texttt{drvr.macos(name)}, \texttt{drvr.android(name)} --- generate
  driver templates
\item
  \texttt{drvr.build(target)}, \texttt{drvr.deploy(target)} ---
  cross-platform build and deploy
\end{itemize}

Comparisons: C++=\texttt{Linux\ module\_init/WDK/WDF} ·
Python=\texttt{ctypes/cffi\ (userspace\ only)} ·
Julia=\texttt{no\ support} · Go=\texttt{TinyGo\ (embedded)} ·
Ruby=\texttt{no\ support} · Rust=\texttt{linux-kernel-module-rust}

\subsubsection{Other Builtins}\label{other-builtins}

\begin{itemize}
\tightlist
\item
  \textbf{\texttt{nlp}} --- Natural Language Processing (delegates to
  \texttt{ai.prompt()}). C++=\texttt{OpenNLP} ·
  Python=\texttt{nltk/spacy} · Julia=\texttt{TextAnalysis.jl} ·
  Rust=\texttt{rust-bert}
\item
  \textbf{\texttt{bev}} --- Bevington data reduction:
  \texttt{data(x,y)}, \texttt{fit\_line()}, \texttt{param(idx)}.
  C++=\texttt{GSL} · Python=\texttt{numpy.polyfit} ·
  Julia=\texttt{GLM.jl} · Rust=\texttt{linregress}
\item
  \textbf{\texttt{file}} --- \texttt{file.write(path,\ content)},
  \texttt{file.exists(path)}. C++=\texttt{std::filesystem} ·
  Python=\texttt{open/os.path} · Julia=\texttt{write/isfile} ·
  Rust=\texttt{std::fs}
\item
  \textbf{\texttt{zip}} --- \texttt{zip.compress(src,\ archive)},
  \texttt{zip.decompress(archive,\ dest)}. C++=\texttt{libzip/miniz} ·
  Python=\texttt{zipfile} · Julia=\texttt{ZipFile.jl} ·
  Rust=\texttt{zip\ crate}
\item
  \textbf{\texttt{img}} ---
  \texttt{img.resize(src,\ width,\ height,\ dest)}.
  C++=\texttt{Magick++} · Python=\texttt{Pillow} ·
  Julia=\texttt{Images.jl} · Rust=\texttt{image\ crate}
\item
  \textbf{\texttt{cpp}} --- \texttt{cpp.repl()}, \texttt{cpp.exec(code)}
  --- C++26 JIT via Cling. Python=\texttt{ctypes/pybind11} ·
  Julia=\texttt{Cxx.jl} · Rust=\texttt{bindgen}
\item
  \textbf{\texttt{itr}} --- \texttt{itr.range(n)},
  \texttt{itr.while(cond\_fun,\ body\_fun)}.
  C++=\texttt{std::ranges::iota} · Python=\texttt{range()} ·
  Julia=\texttt{1:n} · Rust=\texttt{0..n}
\item
  \textbf{\texttt{go}}, \textbf{\texttt{ruby}}, \textbf{\texttt{node}},
  \textbf{\texttt{rust}} --- polyglot \texttt{run(input)} +
  \texttt{go.build()}, \texttt{node.npm\_install()},
  \texttt{rust.cargo\_new()}
\item
  \textbf{\texttt{git}} --- Git VCS: \texttt{git.clone(url,{[}dir{]})}
  shallow clone · \texttt{git.pull()} · \texttt{git.commit({[}msg{]})}
  stage+commit · \texttt{git.push()} · \texttt{git.gitsync()}
  pull+commit+push+show. Alias: \texttt{gitsync}. C++=\texttt{libgit2} ·
  Python=\texttt{gitpython} · Julia=\texttt{LibGit2} ·
  Rust=\texttt{git2\ crate}
\end{itemize}

\begin{center}\rule{0.5\linewidth}{0.5pt}\end{center}

\subsection{Mobile Travel Log App}\label{mobile-travel-log-app}

\begin{Shaded}
\begin{Highlighting}[]
\ExtensionTok{./gem}\NormalTok{ travel\_log.g          }\CommentTok{\# start server on :8080}
\CommentTok{\# Android/iPhone → http://\textless{}host{-}ip\textgreater{}:8080}
\CommentTok{\# macOS/Linux/Win11 → http://localhost:8080}
\end{Highlighting}
\end{Shaded}

Uses the \texttt{mobl} builtin + \texttt{travel\_log.html} (PWA). Voice
dictation → AI NLP → GPS pin → live GeoJSON → Plotly map.

\begin{center}\rule{0.5\linewidth}{0.5pt}\end{center}

\subsection{Getting Started}\label{getting-started}

\subsubsection{Prerequisites}\label{prerequisites}

\begin{itemize}
\tightlist
\item
  \textbf{macOS/Linux}: \texttt{g++} (C++26), \texttt{make},
  \texttt{python3} (3.14+), \texttt{libreadline-dev}. Python 3.14
  headers required for embedded polyglot runtime.
\item
  \textbf{Windows 11}: \href{https://www.msys2.org/}{MSYS2} with
  \texttt{mingw-w64-x86\_64-gcc}, \texttt{make}, and
  \texttt{mingw-w64-x86\_64-python}. Ensure
  \texttt{C:\textbackslash{}msys64\textbackslash{}mingw64\textbackslash{}bin}
  is in your PATH.
\end{itemize}

\subsubsection{Build and Run}\label{build-and-run}

\begin{Shaded}
\begin{Highlighting}[]
\FunctionTok{make}\NormalTok{ all}
\ExtensionTok{./gem}\NormalTok{ tutorial/01\_basics.g}
\end{Highlighting}
\end{Shaded}

On Windows, the binaries will be produced as \texttt{gem.exe} and
\texttt{gem\_test.exe}.

\begin{center}\rule{0.5\linewidth}{0.5pt}\end{center}

\subsection{Windows Support}\label{windows-support}

Gem is now fully compatible with Windows 11 via MSYS2/MinGW-w64. Key
improvements include: - \textbf{Robust Dependency Detection}: Automatic
detection of Python headers and libraries using \texttt{sysconfig}. -
\textbf{Modern API Integration}: Support for \texttt{cpp-httplib} and
other libraries through \texttt{\_WIN32\_WINNT=0x0A00}. -
\textbf{Cross-Platform Signals}: Graceful handling of platform-specific
signals (\texttt{SIGINT}, \texttt{SIGTERM}). - \textbf{Path
Compatibility}: Seamless handling of Windows-style paths and
\texttt{.exe} extensions in compilation mode.

Or enter the REPL:

\begin{Shaded}
\begin{Highlighting}[]
\ExtensionTok{./gem}
\end{Highlighting}
\end{Shaded}

\begin{center}\rule{0.5\linewidth}{0.5pt}\end{center}

\emph{Developed with Gemini CLI and Kiro.}

\end{document}
